% Gustavo Lazzari — Administrador de Sistemas Júnior / Engenheiro de Infraestrutura
% Gerado: 2026-02-26
% Template: Baseado em Rover Resume (https://github.com/subidit/rover-resume)

\documentclass[10pt]{article}

\usepackage[margin=0.6in, a4paper]{geometry}
\setcounter{secnumdepth}{0}
\usepackage{titlesec}
\titlespacing{\section}{0pt}{5pt}{2pt}
\titlespacing{\subsection}{0pt}{3pt}{1pt}
\titlespacing{\subsubsection}{0pt}{2pt}{1pt}
\titleformat{\section}{\large\bfseries\uppercase}{}{}{}[\titlerule]
\titleformat*{\subsubsection}{\normalsize\itshape}
\usepackage{enumitem}
\setlist[itemize]{noitemsep,topsep=1pt,left=0pt .. \parindent}
\setlist[description]{noitemsep,topsep=1pt,itemsep=1pt}
\pagestyle{empty}
\pdfgentounicode=1
\setlength{\parskip}{1pt}
\usepackage[utf8]{inputenc}
\usepackage[T1]{fontenc}
\usepackage{hyperref}
\hypersetup{hidelinks}
\urlstyle{same}

\begin{document}

\begin{center}
	{\LARGE\bfseries Gustavo Lazzari} \\
	\vspace{2pt}
	\href{mailto:gustavotlazzari@gmail.com}{gustavotlazzari@gmail.com} $|$
	+55 41 98773-2772 $|$
	\href{https://www.linkedin.com/in/glazzari}{linkedin.com/in/glazzari} $|$
	\href{https://github.com/GLazzari1428}{github.com/GLazzari1428} $|$
	Curitiba, PR, Brasil
\end{center}

\section{Forma\c{c}\~{a}o Acad\^{e}mica}
\subsection{Pontif\'{i}cia Universidade Cat\'{o}lica do Paran\'{a} $|$ {\normalfont\itshape Bacharelado em Ci\^{e}ncia da Computa\c{c}\~{a}o} \hfill 2024 -- 2028}
\begin{itemize}
	\item Conclus\~{a}o prevista: Junho 2028 $|$ Sistemas Operacionais, Banco de Dados, Seguran\c{c}a de Redes, Estrutura de Dados
\end{itemize}
\subsection{Col\'{e}gio Su\'{i}\c{c}o-Brasileiro de Curitiba $|$ {\normalfont\itshape Ensino M\'{e}dio e Diploma Bil\'{i}ngue IB} \hfill 2017 -- 2020}

\section{Experi\^{e}ncia Profissional}
\subsection{Tribunal de Justi\c{c}a do Paran\'{a} (TJPR) \hfill Curitiba, Brasil}
\subsubsection{Estagi\'{a}rio de Infraestrutura de TI \hfill Ago 2021 -- Dez 2022}
\begin{itemize}
	\item Provisionei e gerenciei infraestrutura de aplica\c{c}\~{o}es em AWS EC2 e GCP Compute Engine, realizando deploy, configura\c{c}\~{a}o e manuten\c{c}\~{a}o de servi\c{c}os em produ\c{c}\~{a}o.
	\item Implementei ambientes de deploy containerizados com Docker e Docker Compose, padronizando workflows de desenvolvimento e produ\c{c}\~{a}o.
	\item Configurei servidores web Nginx e Apache para hospedagem de aplica\c{c}\~{o}es; automatizei pipelines de build e deploy com GitHub Actions.
\end{itemize}

\section{Projetos}
\subsection{Cluster de Servidores Multi-N\'{o} $|$ \normalfont\textit{Infraestrutura Homelab de N\'{i}vel Empresarial} \hfill Em Andamento}
\begin{itemize}
	\item Arquitetei e opero um cluster Proxmox VE de dois n\'{o}s: n\'{o} prim\'{a}rio (Supermicro Xeon E5-2450L, 64GB RAM, Quadro P400) e n\'{o} secund\'{a}rio (Intel i5, 16GB RAM) hospedando 20+ VMs e containers LXC simult\^{a}neos.
	\item Gerencio pools de armazenamento ZFS com experi\^{e}ncia pr\'{a}tica em diagn\'{o}stico de falhas de disco, recupera\c{c}\~{a}o de pools e monitoramento de integridade de dados em discos enterprise.
	\item Projetei estrat\'{e}gia automatizada de backup: n\'{o} secund\'{a}rio orquestra todos os backups de VMs/containers com replica\c{c}\~{a}o offsite para Google Drive para prontid\~{a}o de recupera\c{c}\~{a}o de desastres.
	\item Implantei infraestrutura DNS de alta disponibilidade (PiHole) em ambos os n\'{o}s; sistemas protegidos por nobreak garantindo continuidade de servi\c{c}o.
	\item Configurei PCI passthrough (Jellyfin VM) e mapeamento de dispositivos LXC (Tdarr) para transcodifica\c{c}\~{a}o H.265 acelerada por GPU via Quadro P400.
\end{itemize}

\subsection{Arquitetura de Rede \& Seguran\c{c}a $|$ \normalfont\textit{Firewall, VPN, Zero Trust} \hfill Em Andamento}
\begin{itemize}
	\item Configurei segmenta\c{c}\~{a}o via VLANs (IoT/produ\c{c}\~{a}o) com regras de firewall via OPNsense (migrando para OpenWRT); conhecimento pr\'{a}tico de roteamento, sub-redes e pol\'{i}ticas inter-VLAN.
	\item Implantei proxy reverso Nginx atr\'{a}s de VPN WireGuard para acesso interno seguro; t\'{u}neis Cloudflare com Zero Trust e TOTP para servi\c{c}os p\'{u}blicos.
	\item Implementei gerenciamento automatizado de SSL/TLS via Certbot com desafios DNS Cloudflare para todos os servi\c{c}os expostos.
	\item Auto-hospedo gerenciamento de identidade e credenciais: Vaultwarden (cofre de senhas) e PocketID (provedor de identidade).
\end{itemize}

\section{Habilidades T\'{e}cnicas}
\begin{description}
	\item[Administra\c{c}\~{a}o de Sistemas] Linux (Arch/Debian), Proxmox VE (multi-n\'{o}), Docker, Docker Compose, LXC, systemd, ZFS, RAID
	\item[Redes \& Seguran\c{c}a] OPNsense, OpenWRT, Nginx, WireGuard, Cloudflare (Tunnels/Zero Trust), PiHole, SSL/TLS, VLANs
	\item[Armazenamento \& Backup] Gerenciamento de pools ZFS, recupera\c{c}\~{a}o de falhas de disco, Proxmox Backup, replica\c{c}\~{a}o offsite, monitoramento S.M.A.R.T. (Scrutiny)
	\item[Monitoramento] Prometheus, Grafana, Scrutiny (S.M.A.R.T.)
	\item[Cloud \& Automa\c{c}\~{a}o] AWS (EC2), GCP, GitHub Actions, Bash, Python, YAML/IaC
	\item[Idiomas] Portugu\^{e}s (Nativo), Ingl\^{e}s (Fluente), Alem\~{a}o (Intermedi\'{a}rio)
\end{description}

\section{Certifica\c{c}\~{o}es}
Cambridge Advanced English (CAE -- B2) $|$ Deutsches Sprachdiplom (DSD -- B1/B2) $|$ Em curso: AWS Certified Cloud Practitioner (CLF-C02)

\end{document}
